\documentclass[11pt,a4paper]{article}

\usepackage[utf8]{inputenc}
\usepackage[T1]{fontenc}
\usepackage[english]{babel}
\usepackage{lmodern}
\usepackage{microtype}
\usepackage{amsmath,amssymb,mathtools}
\usepackage{geometry}
\usepackage{xcolor}
\usepackage[hidelinks]{hyperref}

\geometry{left=2.3cm,right=2.3cm,top=2.2cm,bottom=2.3cm}
\setlength{\parindent}{0pt}
\setlength{\parskip}{0.45em}
\allowdisplaybreaks

\definecolor{heading}{RGB}{34,73,110}
\newcommand{\Id}{I}

\begin{document}

\begin{center}
  {\LARGE\bfseries\color{heading}Implicit ABBA Jacobian formulas}\par
  \vspace{0.35em}
  {\large Implicit-function and stage-increment factorizations}
\end{center}

\section{Notation}

Let
\[
 s:=\frac h2,
 \qquad
 W_i:=D_zf(\text{ABBA stage point}_i,\text{stage time}_i),
 \qquad
 S:=W_2+W_3.
\]
Every $W_i$ is a $2\times2$ vector-field Jacobian. Matrix-product order must be
preserved because the $W_i$ do not commute in general.

The implicit multiplier is a function of the physical input,
\[
 \mu_n=\mu_h(z_n),
\]
and its total derivative is denoted by
\[
 \boxed{M:=D\mu_h(z_n).}
\]
The two displaced initial copies therefore satisfy
\[
 Du^{(0)}=\Id+M,
 \qquad
 Dv^{(0)}=\Id-M.
\]

\section{Residual derivatives \texorpdfstring{$K$}{K} and \texorpdfstring{$L$}{L}}

For the projection residual $R(z,\mu)$, define
\[
 K:=D_\mu R,
 \qquad
 L:=D_zR.
\]
In terms of the four stage Jacobians,
\[
 \boxed{
 \begin{aligned}
 K={}&4\Id
 -s(W_1+W_2+W_3+W_4)\\
 &+s^2\bigl(W_4S+SW_1\bigr)
 -s^3W_4SW_1,
 \end{aligned}
 }
\]
and
\[
 \boxed{
 \begin{aligned}
 L={}&s(W_1-W_2-W_3+W_4)\\
 &+s^2\bigl(W_4S-SW_1\bigr)
 +s^3W_4SW_1.
 \end{aligned}
 }
\]
Equivalently, after replacing $S=W_2+W_3$,
\[
 \begin{aligned}
 K={}&4\Id-s(W_1+W_2+W_3+W_4)\\
 &+s^2\bigl[W_4(W_2+W_3)+(W_2+W_3)W_1\bigr]\\
 &-s^3W_4(W_2+W_3)W_1,
 \end{aligned}
\]
\[
 \begin{aligned}
 L={}&s(W_1-W_2-W_3+W_4)\\
 &+s^2\bigl[W_4(W_2+W_3)-(W_2+W_3)W_1\bigr]\\
 &+s^3W_4(W_2+W_3)W_1.
 \end{aligned}
\]

Differentiating the exact-root identity
\[
 R\bigl(z,\mu_h(z)\bigr)=0
\]
gives
\[
 L+KM=0.
\]
If $K$ is invertible, then
\[
 \boxed{M=-K^{-1}L.}
\]
In an implementation this inverse is not formed explicitly: the matrix system
$KM=-L$ is solved instead.

\section{Total derivatives of the four increments}

The physical mean can be written as
\[
 z_{n+1}=z_n+\frac h4(k_1+k_2+k_3+k_4).
\]
The total stage derivatives, including the dependence of $\mu_h$ on $z_n$, obey
\[
 \begin{aligned}
 Dk_1&=W_1(\Id-M),\\
 Du^{(1)}&=\Id+M+sDk_1,\\
 Dk_2&=W_2Du^{(1)},\\
 Dk_3&=W_3Du^{(1)},\\
 Dv^{(2)}&=\Id-M+s(Dk_2+Dk_3),\\
 Dk_4&=W_4Dv^{(2)}.
 \end{aligned}
\]
Eliminating the intermediate derivatives gives the four requested expressions:
\[
 \boxed{Dk_1=W_1(\Id-M),}
\]
\[
 \boxed{
 Dk_2=W_2\bigl[\Id+M+sW_1(\Id-M)\bigr],
 }
\]
\[
 \boxed{
 Dk_3=W_3\bigl[\Id+M+sW_1(\Id-M)\bigr],
 }
\]
and
\[
 \boxed{
 Dk_4=W_4\Bigl\{
 \Id-M+s(W_2+W_3)
 \bigl[\Id+M+sW_1(\Id-M)\bigr]
 \Bigr\}.
 }
\]
Consequently, the stage-increment factorization is
\[
 \boxed{
 J_n^{\mathrm{inc}}
 =\Id+\frac h4(Dk_1+Dk_2+Dk_3+Dk_4).
 }
\]

\section{Equality with the implicit-function factorization}

For completeness, define
\[
 \boxed{
 P=\Id+s(W_1+W_4)+s^2W_4S+s^3W_4SW_1,
 }
\]
\[
 \boxed{
 Q=2\Id-s(W_1+W_4)+s^2W_4S-s^3W_4SW_1.
 }
\]
The derivative of the corrected first copy is
\[
 J_n^{\mathrm{copy}}=P+QM.
\]
Direct symbolic expansion of the stage-increment recurrence gives the identity
\[
 \boxed{
 J_n^{\mathrm{inc}}-(P+QM)=-\frac12(L+KM).
 }
\]
At the implicit root, $L+KM=0$, and hence
\[
 J_n^{\mathrm{inc}}=P+QM.
\]
Finally, substituting $M=-K^{-1}L$ yields
\[
 \boxed{
 J_n^{\mathrm{inc}}
 =\Id+\frac h4\sum_{i=1}^{4}Dk_i
 =P-QK^{-1}L
 =J_n^{\mathrm{IFT}}.
 }
\]

Thus the two calculations are symbolically different derivations and computational
factorizations of the same ideal-root physical Jacobian.

\end{document}
